\section{Discussione}
\paragraph{La scelta dei candidati dipende dalla collezione.}
La Figura~\ref{fig:domains} mostra una regolarità che la media nasconde. Il denso contestuale supera BM25 in ciascuno dei cinque domini di origine StackExchange; BM25 prevale in ciascuno dei cinque domini BRIGHT con dati disponibili. Questa suddivisione segue la provenienza documentata dagli organizzatori \cite{data}. Dimensione delle collezioni, numero di passaggi pertinenti e formulazione delle domande variano insieme alla provenienza. I dati consentono di descrivere l'associazione; attribuirla a una singola proprietà richiederebbe controlli ulteriori. In Robotics, dove il denso ottiene appena 0,0904, la sua fusione con BM25 raggiunge comunque 0,5197. Un retriever debole preso isolatamente può dunque contribuire al primo stadio ibrido. La scelta di un pool congelato è essenziale per stabilire quanto possa aggiungere il riordinamento.

\paragraph{Cosa confrontare con i riferimenti degli organizzatori.}
Per il 2a gli organizzatori pubblicano sul dev completo una nDCG@10 macro di 0,1827 per BM25 sul solo turno e di 0,4379 includendo la storia \cite{data}. La nostra Tabella~\ref{tab:domains} usa una partizione interna di train, copertura parziale dei domini e un'implementazione BM25 nativa Lucene. Una differenza numerica tra le due tabelle confonderebbe sistema e partizione. Il confronto esterno richiede una run su tutti gli 858 turni dev, valutata con gli stessi giudizi e la stessa aggregazione. Un eventuale piazzamento nella classifica del test nascosto richiederà il risultato ufficiale della submission. Per il 2b gli organizzatori pubblicano cinque giudizi distinti a fianco del recupero \cite{evaluation}; un giudizio locale dovrà riportare valutatore e rubrica e mantenere esplicita la sua provenienza.

\paragraph{Quale risultato giustificherebbe il loop?}
Per il 2a una differenza positiva su candidati identici indicherebbe una migliore disposizione delle evidenze accessibili. Esaminare i turni in cui un passaggio entra o esce dalle prime dieci posizioni aiuterebbe a capire se il cambiamento riguarda riferimenti al dialogo o semplici corrispondenze lessicali. Per il 2b una risposta può diventare più completa introducendo un dettaglio privo di sostegno: leggere separatamente completezza e fedeltà evita di perdere questa informazione in un unico punteggio. Un beneficio circoscritto a uno dei compiti sarebbe interpretabile, poiché decidere tra due logit e produrre una sequenza autoregressiva impongono vincoli diversi. A qualità simile, il decoder originale offre il riferimento per giudicare il costo di ogni passaggio ulteriore.

\paragraph{Limiti e riproducibilità.}
I risultati disponibili riguardano un seme, un checkpoint e dieci domini. I soli aggregati di console non permettono di ricostruire intervalli appaiati per conversazione. Per le misure finali occorrerà conservare ranking e risposte per turno, configurazioni effettive, revisioni di codice e pesi, impronte dei candidati, modello di GPU e precisione effettiva. Il profilo destinato a T4 è un vincolo di progetto; i tempi sopra riportati descrivono il singolo log. Il codice del progetto è disponibile all'indirizzo \url{https://github.com/GiovanniIacuzzo/SemEval-2027-Task-1}; la metodologia qui descritta si riferisce alla revisione locale con il loop per entrambi i subtrack.
